\par\noindent\textbf{(1)} 度规 ds²=−(1−2M/r)dt²+(1−2M/r)\textasciicircum{}\{-1\}dr²+r²dφ²，在固有时参数下 E=(1−2M/r)ṫ、L=r²φ̇。利用四速度归一化 u·u=−1 得 ṙ²+V\_eff(r)=E²，V\_eff=(1−2M/r)(1+L²/r²)。
\par\noindent\textbf{(2)} 圆轨道满足 V\_eff'=0。展开 V\_eff=1−2M/r+L²/r²−2ML²/r³，有 V\_eff'=2M/r²−2L²/r³+6ML²/r⁴=0，所以 L²=Mr²/(r−3M)；代入 E²=(r−2M)²/[r(r−3M)]。类时圆轨道要求 r>3M。
\par\noindent\textbf{(3)} 将 L 固定代回 V\_eff''，在圆轨道上其符号与 r−6M 同号。因此 r>6M 稳定，3M<r<6M 不稳定，r\_ISCO=6M；r=6M 为边缘稳定圆轨道。需区别 3M 的光子圆轨道阈值与 6M 的类时最内稳定圆轨道。
\par\noindent\textbf{(4)} 径向轨道方程 u''+u=M/L²+3Mu²。对弱场低偏心率绕Kepler解作微扰，额外 3Mu² 使角频率略减，近心每周进动 Δϖ≈6πM/[a(1−e²)]。还原单位为 6πGM/[c²a(1−e²)]；当 c→∞ 或 a≫GM/c² 时进动消失。
